\chapter{Минимальные гамильтоновы данные multiplicity-среды}
\label{ch:v8-baryon-c0-multiplicity-environment-hamiltonian-minimal-data-gate}

Скалярное gauge-представление не ограничивает оператор энергии среды, а
Real-структура оставляет класс
\begin{equation}
 \mathfrak H_{\rm adm}
 =\{h\in M_3(\mathbb R):h^{\mathsf T}=h\}
 =\operatorname{Sym}_3(\mathbb R),
 \qquad \dim_{\mathbb R}\mathfrak H_{\rm adm}=6.
 \label{eq:v8-baryon-c0-admissible-environment-hamiltonians}
\end{equation}
Если наименьшее собственное значение $h$ невырождено, его спектральный
проектор имеет вид
\begin{equation}
 P=zz^{\mathsf T},
 \qquad z^{\mathsf T}z=1,
 \qquad [z]\in\mathbb{RP}^2.
 \label{eq:v8-baryon-c0-hamiltonian-ground-projector}
\end{equation}
Для задачи выбора направления несущественны положительное изменение
единиц энергии и добавление нуля энергии:
\begin{equation}
 h\sim ah+bI_3,
 \qquad a>0,quad b\in\mathbb R,
 \qquad P_{\min}(ah+bI_3)=P_{\min}(h).
 \label{eq:v8-baryon-c0-ground-projector-affine-equivalence}
\end{equation}

\section{Минимальный одноосный представитель}

Всякое выбранное направление допускает минимальный гамильтониан с
двукратно вырожденным возбуждённым уровнем
\begin{equation}
 h(P;\varepsilon,\Delta)
 =\varepsilon I_3+\Delta(I_3-P),
 \qquad \Delta>0.
 \label{eq:v8-baryon-c0-minimal-axial-environment-hamiltonian}
\end{equation}
Его спектральные данные равны
\begin{equation}
 \Spec h(P;\varepsilon,\Delta)
 =\{\varepsilon^{(1)},(\varepsilon+\Delta)^{(2)}\},
 \qquad P_{\min}(h)=P.
 \label{eq:v8-baryon-c0-minimal-axial-hamiltonian-spectrum}
\end{equation}
Достаточность проверяется без диагонализации:
\begin{equation}
 (h-\varepsilon I_3)P=0,
 \qquad
 (h-\varepsilon I_3)(I_3-P)=\Delta(I_3-P).
 \label{eq:v8-baryon-c0-axial-hamiltonian-projector-identities}
\end{equation}
Любое дополнительное расщепление возбуждённой плоскости не меняет ground
projector и потому избыточно для одной только карты $c_0$.

\section{Конечная температура не даёт точный проектор}

Положим $x=e^{-\beta\Delta}$. Gibbs-состояние минимального представителя
есть
\begin{equation}
 \rho_{\beta,P}
 =\frac{P+x(I_3-P)}{1+2x},
 \qquad 0<x\leq1.
 \label{eq:v8-baryon-c0-axial-hamiltonian-gibbs-state}
\end{equation}
При всяком конечном $\beta\Delta$ выполнено
\begin{equation}
 \rank\rho_{\beta,P}=3,
 \qquad
 \Tr\rho_{\beta,P}^2
 =\frac{1+2x^2}{(1+2x)^2}<1.
 \label{eq:v8-baryon-c0-finite-temperature-full-rank-purity}
\end{equation}
Точный чистый предел имеет вид
\begin{equation}
 \lim_{\beta\Delta\to\infty}\rho_{\beta,P}=P,
 \qquad
 \frac12\lVert\rho_{\beta,P}-P\rVert_1
 =\frac{2x}{1+2x}\longrightarrow0.
 \label{eq:v8-baryon-c0-zero-temperature-projector-limit}
\end{equation}
Следовательно, конечномерный ограниченный гамильтониан при конечной
температуре не создаёт строгое rank-one состояние. Нужен либо предел
$\beta\Delta\to\infty$, либо отдельный протокол чистой подготовки.

\section{Три независимых слоя данных}

Минимальный контракт распадается на
\begin{equation}
 \underbrace{P\in\mathbb{RP}^2}_{\text{направление }c_0},
 \qquad
 \underbrace{\beta\Delta\to\infty}_{\text{точная чистота}},
 \qquad
 \underbrace{\Delta_{\rm phys}>0}_{\text{энергия и время}}.
 \label{eq:v8-baryon-c0-three-layer-environment-data}
\end{equation}
Первый слой необходим для выбора карты. Второй необходим, только если
состояние должно быть получено как Gibbs ground-state limit. Третий не
нужен для самого проектора, но необходим для абсолютной скорости подготовки
и физического времени.

Обратное построение
\begin{equation}
 h_{\rho_*}=\Delta(I_3-\rho_*),
 \qquad \rho_*^2=\rho_*,\quad\rank\rho_*=1
 \label{eq:v8-baryon-c0-projector-to-hamiltonian-circular-construction}
\end{equation}
показывает достаточность любого чистого состояния, но не является его
выводом: искомый селектор уже содержится в $\rho_*$.

После факторизации по $h\sim ah+bI_3$ минимальные данные самого селектора
равны
\begin{equation}
 \mathcal S_{\rm selector}^{\min}=\mathbb{RP}^2,
 \qquad
 N_{\rm derived}(P,\beta\Delta,\Delta_{\rm phys})=0/3.
 \label{eq:v8-baryon-c0-minimal-hamiltonian-data-ledger}
\end{equation}

\begin{theorem}[Минимальность одноосного гамильтониана]
Для выбора одной Real-линии в трёхмерной multiplicity-среде достаточно и,
с точностью до положительного affine-преобразования энергии, необходимо
задать rank-one projector $P\in\mathbb{RP}^2$. Представитель
$h=\varepsilon I_3+\Delta(I_3-P)$ имеет минимальное расщепление $1+2$ и не
содержит лишней ориентации возбуждённой плоскости.
\end{theorem}

\begin{nogocriterion}[Конечный Gibbs-вес не есть чистая подготовка]
Нельзя объявлять $\rho_{\beta,P}$ rank-one состоянием при конечном
$\beta\Delta$: оно имеет ранг три. Нельзя также считать формулу
$h=\Delta(I-P)$ происхождением $P$, поскольку она использует искомый
projector как вход.
\end{nogocriterion}

\begin{protocol}\begin{enumerate}
\item допустимый Hamiltonian-класс: \textbf{$\operatorname{Sym}_3(\mathbb R)$};
\item его вещественная размерность: \textbf{$6$};
\item минимальный спектральный тип: \textbf{$1+2$};
\item селектор после affine-факторизации: \textbf{$P\in\mathbb{RP}^2$};
\item конечный Gibbs-state rank: \textbf{$3$};
\item точная чистота при конечном $\beta\Delta$: \textbf{нет};
\item нулетемпературный предел: \textbf{да, условно};
\item независимых слоёв нового данного: \textbf{$3$};
\item выведено текущим parent: \textbf{$0/3$};
\item одиночный $c_0$ выведен: \textbf{нет};
\item следующий гейт: \textbf{происхождение multiplicity-Hamiltonian из parent}.
\end{enumerate}\end{protocol}

Машинный сертификат сохранён в
\path{s2t_v8_baryon_c0_multiplicity_environment_hamiltonian_minimal_data_gate_results.json}.